% Part: applied-modal-logics
% Chapter: epistemic-logic
% Section: bisimulations

\documentclass[../../../include/open-logic-section]{subfiles}

\begin{document}

\olfileid{aml}{el}{bsd}

\olsection{ద్విసమానుకరణలు}

మన జ్ఞానసంబంధ భాష వ్యక్తీకరణ శక్తి గురించి ఇంకా
ఒక ప్రశ్న మిగిలింది: నమూనాలకూ వాటిలో నిలిచే
సూత్రాలకూ మధ్య సంబంధం ఏమిటి? చట్ర అనురూపత
ఫలితాల నుంచి, ఒక చట్రంలో కొన్ని సూత్రాలు
చెల్లుబాటు అయితే ఆ చట్రానికి కొన్ని ధర్మాలూ
ఉంటాయని చూశాం. అయితే, ఉదాహరణకు, స్వావర్తన
బాణం ఉన్న లోకానికీ, ప్రతి లోకం తదుపరి లోకానికి
దారితీసే అనంత లోకాల గొలుసుకీ మధ్య తేడాను
మన మోడల్ భాష గుర్తించగలదా? అంటే, ఈ రెండు
లోకాల్లో ఒకదాని వద్ద మాత్రమే నిలిచే సూత్రం~$A$
ఏదైనా ఉందా?

సంబంధ నమూనాల మధ్య నిర్మాణపరంగా అవి ఒకేలా
ఉన్నాయని చెప్పడానికి నిర్వచించగల సంబంధమే
ద్విసమానుకరణ. మనం చూడబోయే విధంగా, ఇది
జ్ఞానసంబంధ భాషలో వ్యక్తీకరించగల నమూనాల
మధ్య తుల్యత భావనను పట్టుకుంటుంది.

\begin{defn}[ద్విసమానుకరణ]
$M_1 = \tuple{W_1, R_1, V_1}$, $M_2 = \tuple{W_2, R_2, V_2}$
రెండు సంబంధ నమూనాలు అనుకుందాం. అలాగే
$\mathcal{R} \subseteq W_1 \times W_2$ ఒక ద్విస్థానిక
సంబంధం అనుకుందాం. ప్రతి $\tuple{w_1, w_2} \in \mathcal{R}$కు
కింది షరతులు నెరవేరితే, $\mathcal{R}$ను
\emph{ద్విసమానుకరణ} అంటాం:

\begin{enumerate}
  \item అన్ని \tetoken{ప్రతిజ్ఞావాక్య చరాలు}{!!{propositional variable}s}~$p$కు
  $w_1 \in V_1(p)$ అయితే, అప్పుడే $w_2 \in V_2(p)$.
  \item ప్రతి కర్త $a \in G$, ప్రతి లోకం $v_1 \in W_1$కు,
  $R_{1_a} w_1 v_1$ అయితే, $R_{2_a} w_2 v_2$,
  $\tuple{v_1, v_2} \in
  \mathcal{R}$ అయ్యే ఏదో $v_2 \in W_2$ ఉంటుంది.
  \item ప్రతి కర్త $a \in G$, ప్రతి లోకం $v_2 \in W_2$కు,
   $R_{2_a} w_2 v_2$ అయితే, $R_{1_a} w_1 v_1$,
   $\tuple{v_1, v_2} \in
   \mathcal{R}$ అయ్యే ఏదో $v_1 \in W_1$ ఉంటుంది.
   \sourcecorrection{OLTEAMLELBIS-001}{మూలంలో ఈ రెండు నిబంధనల్లో కర్తల సమితిని $A$గా వ్రాశారు; భాషలో ముందుగా నిర్వచించిన కర్తల సమితి $G$గా సరిచేశాం.}
\end{enumerate}

$M_1$, $M_2$ మధ్య $w_1$, $w_2$ లోకాలను కలిపే
ద్విసమానుకరణ ఉంటే, $\tuple{M_1, w_1} \leftrightarroweq
\tuple{M_2, w_2}$ అని కూడా వ్రాయవచ్చు; అప్పుడు
$\tuple{M_1, w_1}$, $\tuple{M_2, w_2}$లను
\emph{ద్విసమానమైనవి} అంటాం.
\end{defn}

ద్విసమానుకరణ సంబంధంలోని వేర్వేరు నిబంధనలు
వేర్వేరు విషయాలను నిర్ధారిస్తాయి. మొదటి నిబంధన
అన్ని \tetoken{ప్రతిజ్ఞావాక్య చరాలు}{!!{propositional variable}s}పై
ఏకీభావాన్ని నిర్ధారిస్తుంది; అందుచేత ద్విసమాన
లోకాల వద్ద మోడల్-రహిత సూత్రాలు ఒకే విధంగా
సత్యమవుతాయి. వరుసగా ``ముందుకు'', ``వెనక్కి''
అని పిలిచే మిగతా రెండు నిబంధనలు, ప్రాప్యత
సంబంధాల నిర్మాణం ఒకేలా ఉంటుందని నిర్ధారిస్తాయి.

\begin{thm}
$\tuple{M_1, w_1} \leftrightarroweq \tuple{M_2, w_2}$
అయితే, ప్రతి \tetoken{సూత్రం}{!!{formula}}~$!A$కు,
$\mSat{M_1}{!A}[w_1]$ అయితే, అప్పుడే
$\mSat{M_2}{!A}[w_2]$.
\end{thm}

\begin{figure}
  \begin{center}
    \begin{tikzpicture}[modal]
      \node[world] (w1) {$w_1$};
      \node[world] (w2) [above left=of w1]{$w_2$};
      \node[world] (w3) [above right=of w1] {$w_3$};
      \draw[<->] (w1) to node {$a$} (w2);
      \draw[->,loop below] (w1) to node {$a$} (w1);
      \draw[<->] (w1) to [swap] node {$a$} (w3);
      \draw[->,loop above] (w2) to node {$a$} (w2) ;
      \draw[->,loop above] (w3) to node {$a$} (w3) ;

      \node[world](v1)[right of=w1, xshift=1.5in]{$v_1$};
      \node[world](v2)[above of=v1]{$v_2$};
      \draw[<->] (v1) to node {$a$} (v2);
      \draw[->,loop below] (v1) to node {$a$} (v1);
      \draw[->,loop above] (v2) to node {$a$} (v2) ;

      \draw[-,dotted] (w1) to (v1) ;
      \draw[-,dotted,bend left=45] (w2) to (v2) ;
      \draw[-,dotted,bend left=30] (w3) to (v2) ;
    \end{tikzpicture}
  \end{center}
  \caption{రెండు ద్విసమాన నమూనాలు.}
  \ollabel{fig:bisimilar}
\end{figure}

\olref{fig:bisimilar}లో చూపిన రెండు నమూనాలు
ఒకదానితో మరొకటి పూర్తిగా సమానంగా లేకపోయినా,
$w_1$,~$v_1$ లోకాలను కలిపే ద్విసమానుకరణ ఉంది.
అదే ద్విసమానుకరణ $w_2$, $w_3$ రెండింటినీ~$v_2$తో
కలుపుతుంది. ఎందుకంటే వాటి మధ్య తేడాను నిజంగా
గుర్తించగలది మన మోడల్ భాషలో ఏదీ వ్యక్తీకరించలేం.
అయితే $w_2$,~$w_3$ వద్ద వేర్వేరు
\tetoken{ప్రతిజ్ఞావాక్య చరాలు}{!!{propositional variable}s}
సత్యమైతే పరిస్థితి వేరుగా ఉంటుంది.

\end{document}
